\subsubsection{Automation}\label{trien-khai-remediate-automation}

Khắc phục tự động được thực thi qua các playbook Ansible để cấu hình nhanh chóng và nhất quán các quy chuẩn bảo mật trên toàn bộ hệ thống.

\textbf{Bảng chi tiết các hoạt động Remediation tự động:}

{\def\LTcaptype{none} % do not increment counter
\begin{longtable}[]{@{}|
  >{\raggedright\arraybackslash}p{(\linewidth - 8\tabcolsep - 5\arrayrulewidth) * \real{0.0973}}|
  >{\raggedright\arraybackslash}p{(\linewidth - 8\tabcolsep - 5\arrayrulewidth) * \real{0.1936}}|
  >{\raggedright\arraybackslash}p{(\linewidth - 8\tabcolsep - 5\arrayrulewidth) * \real{0.3545}}|
  >{\raggedright\arraybackslash}p{(\linewidth - 8\tabcolsep - 5\arrayrulewidth) * \real{0.3545}}|@{}}
\hline
\begin{minipage}[b]{\linewidth}\raggedright
\textbf{Nhóm section}
\end{minipage} & \begin{minipage}[b]{\linewidth}\raggedright
\textbf{Nội dung Audit Auto}
\end{minipage} & \begin{minipage}[b]{\linewidth}\raggedright
\textbf{Phương thức auto cho tất cả(Module gì?)}
\end{minipage} & \begin{minipage}[b]{\linewidth}\raggedright
\textbf{Kết quả cần đạt/giá trị xuất ra để PASS}
\end{minipage} \\
\hline
\endhead
\multirow{4}{=}{1} & 1.1.3 - 1.1.18: Cấu hình Audit rules tự động. & Module \texttt{ansible.builtin.apt} và \texttt{ansible.builtin.copy}. & Các quy tắc được kích hoạt trên host và dịch vụ \texttt{auditd} khởi chạy thành công. \\
\cline{2-4}
& 1.2.2: Tự động cập nhật Docker. & Module \texttt{ansible.builtin.apt}. & Docker Engine được cập nhật lên phiên bản mới nhất. \\
\hline
2 & & & \\
\hline
3.15-3.24 & Ensure correct permission and file ownership is correct & Dùng module stat để truy xuất quyền folder/file & File ownership và permission phải đúng theo tiêu chí(root:root, quyền 660/644) \\
\hline
4 & & & \\
\hline
\multirow{9}{=}{5} & 5.9, 5.14: Giới hạn port và interface. & \begin{minipage}[t]{\linewidth}\raggedright
\texttt{community.docker.}\\
\texttt{docker\_container} với \texttt{published\_ports}.\strut
\end{minipage} & Port chỉ bind vào interface được phê duyệt. \\
\cline{2-4}
& 5.10, 5.30: Tạo và sử dụng network riêng. & \begin{minipage}[t]{\linewidth}\raggedright
\texttt{community.docker.}\\
\texttt{docker\_network} và\\
\texttt{community.docker.}\\
\texttt{docker\_container}.\strut
\end{minipage} & Container dùng \texttt{secure-net}, không dùng host network hoặc default bridge. \\
\cline{2-4}
& 5.11, 5.12, 5.19, 5.29: Giới hạn tài nguyên. & Thuộc tính \texttt{memory}, \texttt{cpu\_shares}, \texttt{ulimits}, \texttt{pids\_limit}. & Container có giới hạn RAM, CPU priority, file descriptor và PID. \\
\cline{2-4}
& 5.13: Root filesystem read-only. & \texttt{read\_only: true} và các mount \texttt{tmpfs}. & Root filesystem không ghi được; ứng dụng chỉ ghi vào vùng được phép. \\
\cline{2-4}
& 5.15: Giới hạn restart. & \texttt{restart\_policy} và \texttt{restart\_retries}. & \texttt{on-failure} với tối đa 5 lần retry. \\
\cline{2-4}
& 5.26: Chặn leo thang đặc quyền. & \texttt{security\_opts}. & Có \texttt{no-new-privileges}. \\
\cline{2-4}
& 5.27: Thiết lập healthcheck. & Thuộc tính \texttt{healthcheck}. & Healthcheck tồn tại và container đạt trạng thái \texttt{healthy}. \\
\cline{2-4}
& 5.28: Ghim image. & Biến \texttt{image\_name} dùng version tag hoặc digest. & Không dùng \texttt{latest} hoặc image không tag. \\
\cline{2-4}
& 5.4: Thu hẹp capability. & \texttt{cap\_drop} và chỉ add capability cần thiết. & Không giữ capability mặc định không cần thiết. \\
\hline
\multirow{2}{=}{6} & 6.1 - 6.2: Tự động dọn dẹp tài nguyên dư thừa. & Module \texttt{ansible.builtin.cron} hoặc chạy prune tự động. & Giải phóng thành công tài nguyên, đĩa sạch và không tích tụ rác. \\
\hline
\end{longtable}
}

